% apx_cbic-errata.tex -- extracted from chapters/apx_b_errata.tex (round5, author decision 1.1).
% Included with % apx_cbic-errata.tex -- extracted from chapters/apx_b_errata.tex (round5, author decision 1.1).
% Included with % apx_cbic-errata.tex -- extracted from chapters/apx_b_errata.tex (round5, author decision 1.1).
% Included with % apx_cbic-errata.tex -- extracted from chapters/apx_b_errata.tex (round5, author decision 1.1).
% Included with \input{tables/apx_cbic-errata}. Float placement, caption and every value
% are byte-identical to the version that was inline. The extraction was verified by building the
% pre-extraction tree with the same prose edits applied and comparing the RENDERED TEXT of both PDFs
% (pypdfium2 text layer): 260,848 characters each, identical. [correction] These headers previously
% claimed a 'rendered PDF checksum' comparison. No checksum comparison was ever run, and it would
% not have matched: the PDF bytes changed (1,322,758 -> 1,323,852) because this commit also carries
% a prose edit. Caught in review; the claim now names the check that was actually performed.
\begin{table}[htb]
\caption{Content errata of the CBIC 2025 article and the corrections applied in
Chapter~\ref{ch:cbic}. Numeric values quoted in the corrections are the published
table values.}
\label{tab:apx:cbic-errata}
\centering
\small
\begin{tabular}{@{}p{0.42\textwidth}p{0.52\textwidth}@{}}
\toprule
\textbf{Defect in the published article} & \textbf{Correction in the chapter} \\
\midrule
Prose states ``almost four times more wall time''; the convergence table gives
80.88~s for the joint model against a cumulative 34.97~s for the single-task models,
a factor of about 2.3. &
Prose reconciled to the table: ``about 2.3 times the cumulative 34.97~s of the
individual single-task models''. \\
\addlinespace
Prose states the MFLOPs cost of the joint model is ``roughly double''; the
table gives 0.234 MFLOPs for the joint model against 2.315 and 0.012 for the two
single-task models. &
Prose reconciled to the table: the text now states that the table does not show a
higher MFLOPs cost for the joint model and quotes the three cell values; four
follow-on wording edits restrict the overhead claim to wall time. \\
\addlinespace
Broken cross-reference: the label intended for the single-task heads sits on the
Dataset subsection, and the method section points to it as if it described the heads. &
Dataset subsection relabeled; the dangling pointer sentence was removed, since no
section describing the single-task heads exists in the published paper. \\
\addlinespace
Typo ``spatio-tegm mporal''. &
``spatio-temporal''. \\
\addlinespace
Unfilled dataset placeholders ($N_{\text{users}}$, $N_{\text{poi}}$,
$N_{\text{checkins}}$) in the results section. &
Filled with the Florida figures of record for the same Gowalla data, as published in the
CoUrb dataset table: 20{,}301 users, 65{,}009 POIs, and 990{,}518 check-ins. The published
CBIC paper left these three statistics as placeholders; they are taken from the published
CoUrb table so that one corpus figure serves both chapters, and they are the counts that the
extraction pipeline of that time produced, not the counts the current pipeline produces
(Section~\ref{apx:errata:corpus}). \\
\addlinespace
The optimizer subsection states that Nash-MTL's ``reliance on gradient signs rather than
scales inherently balances heterogeneous tasks''. Nash-MTL does not read gradient signs,
and the same subsection states the correct property four paragraphs earlier, listing
scale invariance among the optimizer's axioms. &
Corrected to name the property the method has: its invariance to the scale of the
individual task gradients, the axiom already listed in the same subsection, balances
heterogeneous tasks without manual re-weighting. Checked against the optimizer's own
paper and against the implementation released with the two studies, neither of which
uses gradient signs. The claim is weaker than the published one, and no result depends
on it. \\
\bottomrule
\end{tabular}
\end{table}
. Float placement, caption and every value
% are byte-identical to the version that was inline. The extraction was verified by building the
% pre-extraction tree with the same prose edits applied and comparing the RENDERED TEXT of both PDFs
% (pypdfium2 text layer): 260,848 characters each, identical. [correction] These headers previously
% claimed a 'rendered PDF checksum' comparison. No checksum comparison was ever run, and it would
% not have matched: the PDF bytes changed (1,322,758 -> 1,323,852) because this commit also carries
% a prose edit. Caught in review; the claim now names the check that was actually performed.
\begin{table}[htb]
\caption{Content errata of the CBIC 2025 article and the corrections applied in
Chapter~\ref{ch:cbic}. Numeric values quoted in the corrections are the published
table values.}
\label{tab:apx:cbic-errata}
\centering
\small
\begin{tabular}{@{}p{0.42\textwidth}p{0.52\textwidth}@{}}
\toprule
\textbf{Defect in the published article} & \textbf{Correction in the chapter} \\
\midrule
Prose states ``almost four times more wall time''; the convergence table gives
80.88~s for the joint model against a cumulative 34.97~s for the single-task models,
a factor of about 2.3. &
Prose reconciled to the table: ``about 2.3 times the cumulative 34.97~s of the
individual single-task models''. \\
\addlinespace
Prose states the MFLOPs cost of the joint model is ``roughly double''; the
table gives 0.234 MFLOPs for the joint model against 2.315 and 0.012 for the two
single-task models. &
Prose reconciled to the table: the text now states that the table does not show a
higher MFLOPs cost for the joint model and quotes the three cell values; four
follow-on wording edits restrict the overhead claim to wall time. \\
\addlinespace
Broken cross-reference: the label intended for the single-task heads sits on the
Dataset subsection, and the method section points to it as if it described the heads. &
Dataset subsection relabeled; the dangling pointer sentence was removed, since no
section describing the single-task heads exists in the published paper. \\
\addlinespace
Typo ``spatio-tegm mporal''. &
``spatio-temporal''. \\
\addlinespace
Unfilled dataset placeholders ($N_{\text{users}}$, $N_{\text{poi}}$,
$N_{\text{checkins}}$) in the results section. &
Filled with the Florida figures of record for the same Gowalla data, as published in the
CoUrb dataset table: 20{,}301 users, 65{,}009 POIs, and 990{,}518 check-ins. The published
CBIC paper left these three statistics as placeholders; they are taken from the published
CoUrb table so that one corpus figure serves both chapters, and they are the counts that the
extraction pipeline of that time produced, not the counts the current pipeline produces
(Section~\ref{apx:errata:corpus}). \\
\addlinespace
The optimizer subsection states that Nash-MTL's ``reliance on gradient signs rather than
scales inherently balances heterogeneous tasks''. Nash-MTL does not read gradient signs,
and the same subsection states the correct property four paragraphs earlier, listing
scale invariance among the optimizer's axioms. &
Corrected to name the property the method has: its invariance to the scale of the
individual task gradients, the axiom already listed in the same subsection, balances
heterogeneous tasks without manual re-weighting. Checked against the optimizer's own
paper and against the implementation released with the two studies, neither of which
uses gradient signs. The claim is weaker than the published one, and no result depends
on it. \\
\bottomrule
\end{tabular}
\end{table}
. Float placement, caption and every value
% are byte-identical to the version that was inline. The extraction was verified by building the
% pre-extraction tree with the same prose edits applied and comparing the RENDERED TEXT of both PDFs
% (pypdfium2 text layer): 260,848 characters each, identical. [correction] These headers previously
% claimed a 'rendered PDF checksum' comparison. No checksum comparison was ever run, and it would
% not have matched: the PDF bytes changed (1,322,758 -> 1,323,852) because this commit also carries
% a prose edit. Caught in review; the claim now names the check that was actually performed.
\begin{table}[htb]
\caption{Content errata of the CBIC 2025 article and the corrections applied in
Chapter~\ref{ch:cbic}. Numeric values quoted in the corrections are the published
table values.}
\label{tab:apx:cbic-errata}
\centering
\small
\begin{tabular}{@{}p{0.42\textwidth}p{0.52\textwidth}@{}}
\toprule
\textbf{Defect in the published article} & \textbf{Correction in the chapter} \\
\midrule
Prose states ``almost four times more wall time''; the convergence table gives
80.88~s for the joint model against a cumulative 34.97~s for the single-task models,
a factor of about 2.3. &
Prose reconciled to the table: ``about 2.3 times the cumulative 34.97~s of the
individual single-task models''. \\
\addlinespace
Prose states the MFLOPs cost of the joint model is ``roughly double''; the
table gives 0.234 MFLOPs for the joint model against 2.315 and 0.012 for the two
single-task models. &
Prose reconciled to the table: the text now states that the table does not show a
higher MFLOPs cost for the joint model and quotes the three cell values; four
follow-on wording edits restrict the overhead claim to wall time. \\
\addlinespace
Broken cross-reference: the label intended for the single-task heads sits on the
Dataset subsection, and the method section points to it as if it described the heads. &
Dataset subsection relabeled; the dangling pointer sentence was removed, since no
section describing the single-task heads exists in the published paper. \\
\addlinespace
Typo ``spatio-tegm mporal''. &
``spatio-temporal''. \\
\addlinespace
Unfilled dataset placeholders ($N_{\text{users}}$, $N_{\text{poi}}$,
$N_{\text{checkins}}$) in the results section. &
Filled with the Florida figures of record for the same Gowalla data, as published in the
CoUrb dataset table: 20{,}301 users, 65{,}009 POIs, and 990{,}518 check-ins. The published
CBIC paper left these three statistics as placeholders; they are taken from the published
CoUrb table so that one corpus figure serves both chapters, and they are the counts that the
extraction pipeline of that time produced, not the counts the current pipeline produces
(Section~\ref{apx:errata:corpus}). \\
\addlinespace
The optimizer subsection states that Nash-MTL's ``reliance on gradient signs rather than
scales inherently balances heterogeneous tasks''. Nash-MTL does not read gradient signs,
and the same subsection states the correct property four paragraphs earlier, listing
scale invariance among the optimizer's axioms. &
Corrected to name the property the method has: its invariance to the scale of the
individual task gradients, the axiom already listed in the same subsection, balances
heterogeneous tasks without manual re-weighting. Checked against the optimizer's own
paper and against the implementation released with the two studies, neither of which
uses gradient signs. The claim is weaker than the published one, and no result depends
on it. \\
\bottomrule
\end{tabular}
\end{table}
. Float placement, caption and every value
% are byte-identical to the version that was inline. The extraction was verified by building the
% pre-extraction tree with the same prose edits applied and comparing the RENDERED TEXT of both PDFs
% (pypdfium2 text layer): 260,848 characters each, identical. [correction] These headers previously
% claimed a 'rendered PDF checksum' comparison. No checksum comparison was ever run, and it would
% not have matched: the PDF bytes changed (1,322,758 -> 1,323,852) because this commit also carries
% a prose edit. Caught in review; the claim now names the check that was actually performed.
\begin{table}[htb]
\caption{Content errata of the CBIC 2025 article and the corrections applied in
Chapter~\ref{ch:cbic}. Numeric values quoted in the corrections are the published
table values.}
\label{tab:apx:cbic-errata}
\centering
\small
\begin{tabular}{@{}p{0.42\textwidth}p{0.52\textwidth}@{}}
\toprule
\textbf{Defect in the published article} & \textbf{Correction in the chapter} \\
\midrule
Prose states ``almost four times more wall time''; the convergence table gives
80.88~s for the joint model against a cumulative 34.97~s for the single-task models,
a factor of about 2.3. &
Prose reconciled to the table: ``about 2.3 times the cumulative 34.97~s of the
individual single-task models''. \\
\addlinespace
Prose states the MFLOPs cost of the joint model is ``roughly double''; the
table gives 0.234 MFLOPs for the joint model against 2.315 and 0.012 for the two
single-task models. &
Prose reconciled to the table: the text now states that the table does not show a
higher MFLOPs cost for the joint model and quotes the three cell values; four
follow-on wording edits restrict the overhead claim to wall time. \\
\addlinespace
Broken cross-reference: the label intended for the single-task heads sits on the
Dataset subsection, and the method section points to it as if it described the heads. &
Dataset subsection relabeled; the dangling pointer sentence was removed, since no
section describing the single-task heads exists in the published paper. \\
\addlinespace
Typo ``spatio-tegm mporal''. &
``spatio-temporal''. \\
\addlinespace
Unfilled dataset placeholders ($N_{\text{users}}$, $N_{\text{poi}}$,
$N_{\text{checkins}}$) in the results section. &
Filled with the Florida figures of record for the same Gowalla data, as published in the
CoUrb dataset table: 20{,}301 users, 65{,}009 POIs, and 990{,}518 check-ins. The published
CBIC paper left these three statistics as placeholders; they are taken from the published
CoUrb table so that one corpus figure serves both chapters, and they are the counts that the
extraction pipeline of that time produced, not the counts the current pipeline produces
(Section~\ref{apx:errata:corpus}). \\
\addlinespace
The optimizer subsection states that Nash-MTL's ``reliance on gradient signs rather than
scales inherently balances heterogeneous tasks''. Nash-MTL does not read gradient signs,
and the same subsection states the correct property four paragraphs earlier, listing
scale invariance among the optimizer's axioms. &
Corrected to name the property the method has: its invariance to the scale of the
individual task gradients, the axiom already listed in the same subsection, balances
heterogeneous tasks without manual re-weighting. Checked against the optimizer's own
paper and against the implementation released with the two studies, neither of which
uses gradient signs. The claim is weaker than the published one, and no result depends
on it. \\
\bottomrule
\end{tabular}
\end{table}
